% Input after the angular cone formula (AG.2).
\subsection{小齐次次数的完整弱型渐近}\label{subsec:small-homogeneity-weak}

沿用 (AG.1)--(AG.2) 的范数球 $K$、$\rho=\|\cdot\|_K$、$V=|K|$ 与角向概率 $\sigma_K$。
固定可测角向函数 $0<a<\infty$（$\sigma_K$-a.e.），将它零次齐次延拓，并假设
\[
0<E_a:=\int a\,d\sigma_K=\frac1V\int_Ka(x)\,dx<\infty,
\qquad a_*:=\operatorname*{ess\,inf}_{\sigma_K}a.
\]
对任意正局部可积权 $v$，记经典 $A_1$ 特征量\cite{Muckenhoupt1972}为
\[
[v]_{A_1(K)}:=\sup_{Q=x+rK}
\frac{|Q|^{-1}\int_Qv}{\operatorname*{ess\,inf}_Qv}. \tag{SA.1}
\]
分母为零时该比取 $\infty$。本节的弱型条件与近单位首项的 (H.5a) 分开核验；二阶公式仍使用 (H.1) 及相应范围。
对 $0<\tau<n$，置 $d\mu_\tau=h_\tau dx$，其中
$h_\tau(x)=\rho(x)^{\tau-n}a(x)$。令 $M_K$ 是 Lebesgue 平均的中心 $K$-球极大算子，且
\[
\mathcal W_\tau:=W^{\rm out}_{K,h_\tau}=\sup_{\substack{f\ge0,\ f\in L^1_{\rm loc}\\0<\int f h_\tau<\infty}}
\frac{\mu_\tau(\{M_Kf>1\})}{\int f h_\tau}.
\]
记 $B_K$ 为 (G.1) 中的测度一致中心弱型上界；一般可取 $3^n$，方体可取 $2^n$。


\begin{proposition}[齐次权的有限性判别]\label{prop:angular-a1-classification}
以下条件等价：$\mathcal W_{\tau_0}<\infty$ 对某个 $0<\tau_0<n$ 成立；
$a\in A_1(K)$；以及对每个 $0<\tau<n$ 均有 $\mathcal W_\tau<\infty$。
它们也等价于 $h_\tau\in A_1(K)$ 对某个允许的 $\tau$ 成立。条件成立时 $a_*>0$。
\end{proposition}
\begin{proof}
这是经典 $A_1$ 子集检验在当前齐次族中的应用，下面给出所需版本。
若 $\mathcal W_\tau<\infty$，取正测度子集 $E\subset Q=x+rK$。
对所有 $z\in Q$ 有 $Q\subset z+2rK$，故
$M_K\mathbf1_E(z)\ge |E|/(2^n|Q|)$。以略小的严格阈值应用弱型界，再从下逼近，得到
\[
|Q|^{-1}\int_Qh_\tau\le2^n\mathcal W_\tau |E|^{-1}\int_Eh_\tau.
\]
令 $E$ 逼近 $Q$ 中密度的本质下确界，即得
$[h_\tau]_{A_1(K)}\le2^n\mathcal W_\tau$。
反之，逐球比较给出
\[M_Kf\le[h_\tau]_{A_1(K)}M_{K,\mu_\tau}f.\]
结合 (G.1)，得到有限弱型界。

又 $\operatorname*{ess\,inf}_Kh_\tau=a_*$：下界由 $\rho^{\tau-n}\ge1$，上界由靠近
$\partial K$ 的低角正测度锥环带得到。故当 $[h_\tau]_{A_1(K)}<\infty$ 时，
$a_*>0$ 且 $nE_a/(\tau a_*)\le[h_\tau]_{A_1(K)}$。
若 $r<\rho(x)/2$，写 $a=\rho^{n-\tau}h_\tau$，可得
$a_Q/\operatorname*{ess\,inf}_Qa\le3^{n-\tau}[h_\tau]_{A_1(K)}$。
若 $r\ge\rho(x)/2$，则 $Q\subset3rK$，该比至多 $3^nE_a/a_*$。
于是 $a\in A_1(K)$。

反过来，令 $A_a=[a]_{A_1(K)}<\infty$。取 $Q=K$ 得 $E_a/a_*\le A_a$。
对上述小球，$h_\tau$ 的平均与下确界之比至多 $3^{n-\tau}A_a$；对大球，该比至多
\[
\frac{nE_a}{\tau a_*}\left(\frac{\rho(x)+r}{r}\right)^n
\le\frac{nE_a}{\tau a_*}3^n.
\]
所以 $[h_\tau]_{A_1(K)}\le3^n(n/\tau)A_a$ 对每个允许的 $\tau$ 都有限。
\end{proof}

以下假设 $A_a=[a]_{A_1(K)}<\infty$。固定 $q>1$、$0<\eta<1$，记
\[
N=N_q(\eta):=\left\lceil\frac{\log(2/\eta)}{\log q}\right\rceil,
\qquad s_j=q^j/2\quad(0\le j<N),\qquad
\Xi_{n,\tau}(q,\eta):=\sum_{j=0}^{N-1}\frac{(1+qs_j)^{n-\tau}}{s_j^n}.
\]
\begin{theorem}[三尺度上下界]\label{thm:small-homogeneity-weak}
在上述记号下，有
\begin{align}
\mathcal W_\tau\le{}&B_K A_a3^{n-\tau}
+\frac{nE_a}{a_*(n-\tau)}\Xi_{n,\tau}(q,\eta)
+\frac{nE_a}{a_*\tau}\eta^\tau(1+\eta)^{n-\tau}. \tag{SW.1}\label{eq:sw-upper}
\end{align}
并且
\begin{equation}
\mathcal W_\tau\ge\frac{nE_a}{a_*\tau}\,
\frac{\tau^\tau(n-\tau)^{n-\tau}}{n^n}. \tag{SW.2}\label{eq:sw-lower}
\end{equation}
故
\[\boxed{\displaystyle\lim_{\tau\downarrow0}\tau \mathcal W_\tau=\frac{nE_a}{a_*}.}\tag{SW.3}\label{eq:sw-limit}\]
\end{theorem}

\begin{proof}
固定满足条件的 $f$，写 $A=\int f h_\tau$ 与 $H=\{M_Kf>1\}$。输入自动局部 Lebesgue 可积，因为 $h_\tau$ 在每个有界集有正的本质下界。
对 $x\ne0$，将严格见证 $x+rK$（即 $\int_{x+rK}f>Vr^n$）按
\[
r<\rho(x)/2,\qquad \rho(x)/2<r<\rho(x)/\eta,\qquad r>\rho(x)/\eta \tag{SW.4}
\]
分类，令 $H_j$ 为存在第 $j$ 类见证的点集，允许重叠。阈值处的见证可微扰半径而保持严格平均；
各类因而覆盖 $H$ 除原点外的部分。固定半径平均连续，故以有理半径刻画可见各类可测；原点为零测集。

若 $x\in H_0$，径向因子在 $Q=x+rK$ 上的比至多 $3^{n-\tau}$，而
$a_Q\le A_a\operatorname*{ess\,inf}_Qa$。故
\[
\frac{\int_Qf\,d\mu_\tau}{\mu_\tau(Q)}>A_a^{-1}3^{\tau-n},
\qquad
\mu_\tau(H_0)\le B_K A_a3^{n-\tau}A. \tag{SW.5}\label{eq:sw-local}
\]

将中尺度见证进一步按 $s_j\rho(x)\le r<qs_j\rho(x)$ 分层。
因 $q^N/2\ge\eta^{-1}$，这些层覆盖所需半径比；边界仍按严格平均的连续性微扰处理。
第 $j$ 层上 $x+rK\subset(1+qs_j)\rho(x)K$ 且 $\int_{x+rK}f>Vs_j^n\rho(x)^n$，所以
\[
T_j(x):=\frac1{Vs_j^n\rho(x)^n}\int_{(1+qs_j)\rho(x)K}f(y)\,dy>1.
\]
Markov、Tonelli 与 (AG.2) 给出
\begin{align*}
\mu_\tau(H_1)
&\le\sum_{j=0}^{N-1}\int T_j\,d\mu_\tau\\
&=\sum_{j=0}^{N-1}\int f(y)\frac1{Vs_j^n}
\int_{\rho(x)\ge\rho(y)/(1+qs_j)}\rho(x)^{\tau-2n}a(x)\,dx\,dy\\
&\le\frac{nE_a}{a_*(n-\tau)}\Xi_{n,\tau}(q,\eta)A. \tag{SW.6}\label{eq:sw-middle}
\end{align*}
内层积分的径向幂为 $s^{\tau-n-1}ds$。逐层计算保留见证半径比，避免把所有中尺度见证放入同一个大球。

若 $x\in H_2$，则 $x+rK\subset(r+\rho(x))K$ 且 $\rho(x)<\eta r$。因 $\tau-n<0$，
\[
A\ge a_*(r+\rho(x))^{\tau-n}\int_{x+rK}f
>a_*Vr^\tau(1+\eta)^{\tau-n},
\]
于是 $\rho(x)^\tau<\eta^\tau(1+\eta)^{n-\tau}A/(a_*V)$。而
$\mu_\tau(sK)=nVE_as^\tau/\tau$，从而
\[\mu_\tau(H_2)\le\frac{nE_a}{a_*\tau}\eta^\tau(1+\eta)^{n-\tau}A.\tag{SW.7}\label{eq:sw-central}\]
合并 \eqref{eq:sw-local}--\eqref{eq:sw-central} 即得 \eqref{eq:sw-upper}。

为证下界，固定 $\epsilon>0$。低角集合 $\{a<a_*+\epsilon\}$ 在任意固定环带有正测度，
故其中可取 $h_\tau$ 的非零 Lebesgue 点 $z$，满足 $a(z)\le a_*+\epsilon$。令 $s=\rho(z)$ 且 $r>1$。
原子 $Vr^ns^n\delta_z$ 的严格光晕包含 $(r-1)s\operatorname{int}K$：对其内每个 $x$，
$\rho(z-x)<rs$，故存在半径小于 $rs$ 而仍含 $z$ 的中心见证球。用 $z$ 附近小 $K$-球上的
常数质量输入逼近该原子；Lebesgue 点性质令成本趋于 $Vr^ns^\tau a(z)$，且每个紧包含于上述
开球的集合最终仍在严格光晕中。穷尽该开球，得到
\[
\mathcal W_\tau\ge\frac{nE_a}{\tau(a_*+\epsilon)}\frac{(r-1)^\tau}{r^n}.
\]
令 $\epsilon\downarrow0$，并取 $r=n/(n-\tau)$，即得 \eqref{eq:sw-lower}。

固定 $\eta$，把 \eqref{eq:sw-upper} 乘以 $\tau$ 后令 $\tau\downarrow0$，有
\[\limsup_{\tau\downarrow0}\tau \mathcal W_\tau\le\frac{nE_a}{a_*}(1+\eta)^n.\]
再令 $\eta\downarrow0$，并和 \eqref{eq:sw-lower} 合并，即得 \eqref{eq:sw-limit}。
\end{proof}

\begin{corollary}[一维幂权]\label{cor:small-homogeneity-one-dimensional}
当 $n=1$、$\rho(x)=|x|$、$a\equiv1$ 时，
\[
W_\tau\le6+\frac{4N_2(\eta)}{1-\tau}
+\frac{\eta^\tau(1+\eta)^{1-\tau}}{\tau},
\qquad\lim_{\tau\downarrow0}\tau W_\tau=1.
\]
\end{corollary}


\begin{corollary}[一致对数余项]\label{cor:small-homogeneity-uniform}
固定 $n\ge1$、$R\ge1$，并令 $L=nE_a/a_*$。对全部范数球 $K$ 及
$[a]_{A_1(K)}\le R$ 的零次齐次轮廓，$0<\tau<\min\{1,n/2\}$ 时有
\[
1-\tau[1+\log(n/\tau)]\le\frac{\tau\mathcal W_\tau}{L}
\le1+\tau\left[\frac{9^nR}{n}+2^n-1+
\frac{4^nN_2(\tau)}{n-\tau}\right]. \tag{SW.8}
\]
特别地，相对误差为 $O_{n,R}(\tau\log(1/\tau))$。进一步一致有
\[
1-\tau\log(1/\tau)-O_n(\tau)
\le\frac{\tau\mathcal W_\tau}{L}
\le1+(e-1)\tau\log(1/\tau)+O_{n,R}(\tau). \tag{SW.8a}
\]
这些是双侧界，不给出真实第二项或其精确系数。
\end{corollary}
\begin{proof}
下界因子 $L_0=\tau^\tau(n-\tau)^{n-\tau}/n^n$ 满足
$1-L_0\le-\log L_0\le\tau[1+\log(n/\tau)]$，并有
$L_0=1-\tau\log(1/\tau)+O_n(\tau)$。
取 $q=2$、$\eta=\tau$。因 $s_j\ge1/2$，
$(1+2s_j)^{n-\tau}/s_j^n\le4^n$。将 (SW.1) 除以 $L/\tau$，
局部项至多 $9^nR\tau/n$，最后一项不超过
$(1+\tau)^n\le1+(2^n-1)\tau$，遂得 (SW.8)。

为改进对数项常数，固定 $q>1$、$\eta=\tau$。由 $s_j^{-\tau}\le2^\tau$ 及二项式公式，
\[
\Xi_{n,\tau}(q,\tau)\le
2^\tau\sum_{j=0}^{N-1}(q+s_j^{-1})^n
\le2^\tau(q^nN+C_{n,q}),
\]
其中 $C_{n,q}=\sum_{k=1}^n\binom nk q^{n-k}2^k/(1-q^{-k})$，与 $N,K,a$ 无关。
又 $N=\log(1/\tau)/\log q+O_q(1)$，而
$\tau^\tau(1+\tau)^{n-\tau}=1-\tau\log(1/\tau)+O_n(\tau)$。
代入 (SW.1) 得上界的对数项系数 $q^n/(n\log q)-1$。
取 $q=e^{1/n}$ 即为 $e-1$。局部项仍为 $O_{n,R}(\tau)$，证明了一致性。
\end{proof}

系数 $L\ge n$，等号恰在 $a$ 几乎处处为常数时成立。故与平均球匹配的径向幂权
也恰好最小化这一领先系数。另一方面，若进一步 $a$ 有界，则对固定 $a$，(AG.7) 给出的近单位水平系数
$S_\tau=\|a\|_{n/(n-\tau)}/E_a\to1$，由控制收敛直接得到；
小齐次次数的临界行为仍通过 $E_a/a_*$ 保留角向信息。

对固定的正乘积指数向量 $\boldsymbol b$，(OW.2) 进一步给出有确定极限顺序的结论
\[
\lim_{\tau\downarrow0}\tau
\left[\lim_{\lambda\downarrow\lambda_c}(\lambda-\lambda_c)
C^{\boldsymbol b,\rm out}_{K,h_\tau}(\lambda)\right]
=\rho_c\lambda_c\frac{nE_a}{a_*}. \tag{SW.9}
\]
这里先对每个固定权取临界上侧极限，再令 $\tau\downarrow0$。
特别地，一维径向权、等指数双线性模型的临界留数渐近为 $1/(2\tau)$。

\begin{proposition}[立方体上的显式无界 \texorpdfstring{$A_1$}{A1} 剖面]\label{prop:unbounded-angular-cube}
取 $\rho(x)=\|x\|_\infty$，$K=[-1,1]^n$，并令 $0<\gamma<1$，
\[
 a_\gamma(x)=\left(\frac{\|x\|_\infty}{|x_1|}\right)^\gamma\quad(x_1\ne0).
\]
超平面 $x_1=0$ 上任意定义。则 $a_\gamma\in A_1(K)$、$a_*=1$，并且
\[
 E_{a_\gamma}=\frac{n-\gamma}{n(1-\gamma)},\qquad
 [a_\gamma]_{A_1(K)}\le
 \max\left\{\frac{2\,3^\gamma}{1-\gamma},
 \frac{3^n(n-\gamma)}{n(1-\gamma)}\right\}.\tag{SA.2}\label{eq:angular-cube-a1}
\]
故 $\lim_{\tau\downarrow0}\tau \mathcal W_\tau=(n-\gamma)/(1-\gamma)$。当 $n\ge2$ 时该剖面真正无界。
\end{proposition}

\begin{proof}
$a_\gamma\ge1$，且 $|x_1|=\|x\|_\infty$ 的集合有正测度，故 $a_*=1$。对 $n\ge2$，令
$t=|x_1|$、$m=\max_{2\le j\le n}|x_j|$，在单位立方体上先对 $m$ 积分，得
\[
 E_{a_\gamma}=\int_0^1\left[t^{n-1}+\frac{n-1}{n-1+\gamma}
 (t^{-\gamma}-t^{n-1})\right]dt=\frac{n-\gamma}{n(1-\gamma)}.
\]
($n=1$ 时 $a_\gamma\equiv1$。) 对局部立方体 $Q=x+rK$，$r<\|x\|_\infty/2$，一维
$w(t)=|t|^{-\gamma}$ 满足
\[
 \frac{|I|^{-1}\int_Iw}{\mathop{\rm ess\,inf}_Iw}\le\frac2{1-\gamma}.
\]
将其用于第一坐标投影，再用 $\|y\|_\infty^\gamma$ 在 $Q$ 上的比至多 $3^\gamma$，得到
$2\,3^\gamma/(1-\gamma)$。同号区间的比为
$(1-q^{1-\gamma})/((1-\gamma)(1-q))$，跨零区间的比为
$(1+q^{1-\gamma})/((1-\gamma)(1+q))$，故上述界成立。若
$r\ge\|x\|_\infty/2$，则 $Q\subset3rK$；齐次性和 $a_\gamma\ge1$ 给出该比至多
$3^nE_{a_\gamma}$。这证明 \eqref{eq:angular-cube-a1}，定理给出极限。
\end{proof}

\begin{remark}[为何不能只假设 $E_a<\infty$ 与 $a_*>0$]\label{rem:angular-spike-failure}
在二维立方体中令
\[
 a=1+\sum_{k\ge1}2^k\mathbf1_{\{x_2>0,\ x_1/x_2\in I_k\}},
 \qquad I_k=(2^{-k},2^{-k}+2^{-3k-4}).
\]
则 $a_*=1$ 而 $E_a=1+\sum_{k\ge1}2^k|I_k|/8=385/384$。但在
$z_k=(2^{-k},1)$ 的充分小对称立方体上，第 $k$ 个峰在穿过中心的直线一侧占一半，其他峰可避开，
故其 $A_1$ 比至少 $1+2^{k-1}$。于是 $a\notin A_1(K)$，定理说明每个 $\mathcal W_\tau$ 都无穷。
因此有限角向均值和正下界仍不足以保证完整弱型有限性。
\end{remark}
